% =====================================================================
% FORMULATIONS AND METHODOLOGY -- layers, equations, proposed metrics,
% search.  Deliberately orthogonal to the experimental configuration:
% simulator, workloads/traffic and the measurement protocol live in
% §Results.  Moved out of here: substrate block, the "why 3D" paragraph,
% workloads/traffic model, measurement protocol.
% Metric definitions currently duplicated at results.tex §V.B -- delete
% them there if this section is adopted.
% Budget: ~0.8 pp.
% =====================================================================
\section{Formulations and Methodology}
\label{sec:method}
DNN traffic on a tiled IMC substrate is not the traffic classical NoC design
assumes. It is a funnel rather than a permutation: roughly 82\% of a block's NoC
bytes are partial-sum reduction converging on a handful of accumulator tiles,
whose single ejection port is 18--148$\times$ oversubscribed at true model rates.
The binding constraint is therefore a port rather than a link, and it is
invariant to placement and routing --- every placement reachable by ordinary
means leaves the peak link load below that port floor. The traffic graph is
itself a design variable, not an input: how much reduction is internalised
on-tile is decided by the packing, which moves the port floor $5.8\times$ before
any placement exists. It is also non-stationary --- layer-sequential execution
gives each flow a disjoint activity window inside one inference period, so peak
and average diverge by $1.7$--$9.9\times$, separating feasibility from
performance --- and its congestion is transient and clustered at phase
boundaries, with drain tails of 0.8--1.5\,k cycles against phase windows of
5--23\,k, precisely the timescale an adaptive policy's cost field must track.
Finally, the endpoint that matters is the tail rather than the mean, because
reduction is a barrier in which the accumulator waits for the slowest arriving
partial sum; below saturation, throughput is uninformative --- across
configurations spanning $260\%$ in delay, delivered throughput varied by
$0.83\%$ [REF].

This section defines the quantities that characterisation makes necessary, the
design layers that set them, and the search that optimises them.

\subsection{The three design layers}
\label{sec:layers}
The top layer is \emph{architecture and packing}. For a given hardware design it fixes how many crossbars share a tile (packing density, PD), and it decides the \emph{layer-to-tile} mapping, i.e.\ how a DNN layer's rows (reduce, $r$) and columns (scatter, $s$) are grouped across those crossbars (packing orientation, PO). Together, PD and PO fix how much of a layer's traffic has to pass through a single tile port --- a floor on delay that no later decision can lower. The second layer is \emph{tile-to-node} placement, or \emph{mapping}, which assigns NoC mesh nodes to those tiles and sets the peak single-link load across the NoC. The third is the \emph{routing and selection} policy, which decides, at run time, how much of that load can be steered off that link.

The layers are searched in that order, and the search is nested rather than
sequential: each PO induces its own traffic graph \emph{and} its own placement
space, so two orientations are comparable only when each is evaluated at its
own best placement. Comparing a tuned placement of one orientation against an
arbitrary placement of another confounds two layers.

\subsection{Packing: PD, PO and the port floor}
\label{sec:packing}
A packing $(c,r,s)$ with $r\cdot s=c$ assigns the crossbars of a DNN layer to
tiles: $r$ groups crossbars along the reduction dimension, internalising
partial-sum traffic on-tile, and $s$ groups them along the scatter dimension,
internalising input broadcast. It induces a \emph{traffic graph} whose nodes
are logical tiles and whose arcs are flows. Flow $f$ carries rate $\lambda_f$
flits/cycle from tile $u_f$ to tile $v_f$ and is active over the window
$[t^{\mathrm{on}}_f, t^{\mathrm{off}}_f)$ of the inference period $T$.

Because injection and ejection are independent channels, a tile's two port
rates bind separately. Over the elementary intervals $\mathcal{I}$ --- the
maximal windows on which the active flow set is constant, cut at every
$t^{\mathrm{on}}_f$ and $t^{\mathrm{off}}_f$ --- the peak injection and
ejection loads and the \emph{port floor} are
\begin{align}
  \mathrm{PIL} &= \max_{I\in\mathcal{I}}\ \max_{n}
    \sum_{f\,:\,u_f=n,\ f\ \mathrm{active}\ \mathrm{in}\ I} \lambda_f,
  \label{eq:pil}\\
  \mathrm{PEL} &= \max_{I\in\mathcal{I}}\ \max_{n}
    \sum_{f\,:\,v_f=n,\ f\ \mathrm{active}\ \mathrm{in}\ I} \lambda_f,
  \label{eq:pel}\\
  \PF &= \max(\mathrm{PIL},\mathrm{PEL}).
  \label{eq:pf}
\end{align}
Equations~\eqref{eq:pil}--\eqref{eq:pf} involve no placement variable, so
$\PF$ is \emph{placement-invariant}: remapping relabels which physical node is
hot, never how hot it is. It is the quantity the packing layer alone controls,
and $1/\PF$ upper-bounds the load the design can sustain in a burst.

\subsection{Placement: peak link load}
\label{sec:placement}
A placement is an injective map $\pi$ from tiles to mesh nodes. Under a
routing function $\mathcal{R}$, let $\varphi_f(\ell,\pi)\in[0,1]$ be the
fraction of flow $f$'s minimal $\mathcal{R}$-admissible paths that traverse
directed link $\ell$, so that an even spread over admissible paths gives the
offered link load $\Lambda_\ell = \sum_f \lambda_f\,\varphi_f(\ell,\pi)$,
again maximised over $\mathcal{I}$. The placement layer is then characterised
by
\begin{equation}
  \PL = \max_{\ell} \Lambda_\ell,
  \qquad
  \mathrm{CC} = \sum_f \lambda_f\, h\!\left(\pi(u_f),\pi(v_f)\right),
  \label{eq:pl}
\end{equation}
with $h$ the Manhattan hop distance. $\PL$ is the only term a placement moves;
$\mathrm{CC}$ is the hop-count energy proxy. Writing $\ell^{*}$ for the
argmax link, the two layers meet at the \emph{binding load}
\begin{equation}
  \BIND = \max(\PF, \PL),
  \label{eq:bind}
\end{equation}
which is the predictor the flow optimises: below $\PF$ the placement term is
inactive and delay is flat in $\PL$; above it $\PL$ binds and delay is steep.

\subsection{Routing and selection}
\label{sec:routing}
Routing is odd-even-balanced adaptive routing~\cite{dahir2013oeb}, a
deadlock-free 3D turn model that admits, for most source--destination pairs,
several minimal paths; $\mathcal{R}$ in \eqref{eq:pl} is this admissible set.
Which admissible output a packet takes is the \emph{selection} decision, and
we compare two policies. \textbf{BL} is local: it picks the admissible output
whose downstream input buffer has the most free slots, so its view of
congestion extends one hop. \textbf{DP} is distributed and multi-hop: each
router maintains, per destination, a congestion cost-to-go estimate obtained
by relaxing its neighbours' estimates over admissible turns, and selects the
output minimising local occupancy plus that neighbour's cost-to-go. DP
therefore steers around congestion several hops ahead, at the cost of a field
that must reconverge whenever the traffic pattern changes.

\subsection{Proposed metrics}
\label{sec:metrics}
Selection can only divert load that has somewhere to go. For link $\ell$, the
share of its load carried by flows with an alternative is the
\emph{escapable fraction}
\begin{equation}
  E(\ell) = 1 - \frac{1}{\Lambda_\ell}
            \sum_{f\,:\,\varphi_f(\ell)=1} \lambda_f ,
  \label{eq:e}
\end{equation}
the complement of the \emph{forced} load, i.e.\ of the flows every one of whose
admissible paths crosses $\ell$. We write $E = E(\ell^{*})$.

A single hot link, however, does not describe how the escape room is
distributed across the links that actually congest. Let $\mathcal{L}_q$ be the
top-$q$ fraction of links ordered by $\Lambda_\ell$, and let
$E_q = \sum_{\ell\in\mathcal{L}_q} E(\ell)\Lambda_\ell \big/
       \sum_{\ell\in\mathcal{L}_q}\Lambda_\ell$
be the load-weighted escapable fraction of that tier. We propose the
\emph{escape slope}
\begin{equation}
  \mathrm{ES} = E_{20} - E_{5},
  \label{eq:es}
\end{equation}
the change in escape room between the top 5\% and the top 20\% of links.
$\mathrm{ES}$ is computed offline from the traffic graph and the permutation
alone, is invariant to the load scale, and --- unlike $E$ --- is free at fixed
$\PL$ and $\mathrm{CC}$, which is what makes it usable as a refinement
objective rather than a description.

Finally, because different placements saturate at different loads, capacity is
measured per mapping: $\kstar$ is the injection scale at which a mapping
reaches twice its \emph{own} free-flow delay, interpolated from a load ladder.
Mappings are compared at their own $\kstar$, never at one global load.

\subsection{The design flow and its search}
\label{sec:search}
The flow instantiates the layer order above: choose the orientation that
minimises $\PF$; place at minimum $\mathrm{CC}$; refine that placement by
maximising $\mathrm{ES}$ subject to $\mathrm{CC}\le\mathrm{CC}_{\min}$ and
$\PL\le\PL(\mathrm{CC}_{\min})$, which costs nothing at either upstream layer;
minimise $\PL$ if further search is affordable; and run adaptive selection
throughout.

All of \eqref{eq:pil}--\eqref{eq:es} are evaluated from the traffic graph and
the permutation alone, with no simulation. Placements are searched by a
guarded hill-climb seeded from a pool of $1\,000$ random permutations: a step
draws 60 candidate moves --- a transposition of two occupied nodes, or, with
probability $0.3$, a relocation of one tile to an idle node --- discards any
candidate violating the active constraint, and accepts the best strictly
improving one. Applying constraints inside the move filter rather than as a
post-hoc rejection is what makes the guarded objectives above searchable
directly. All searches are seeded and reproducible.
